(1) 在相互作用绘景，\(\rho\) 满足
\[
\dot\rho=-\frac{i\Omega_R}{2}[\sigma_x,\rho]
+\gamma\left(\sigma_-\rho\sigma_+
-\frac12\{\sigma_+\sigma_-,\rho\}\right).
\]
由定义 \(v=2\Im\rho_{eg},\,w=\rho_{ee}-\rho_{gg}\) 计算
\[
\dot v=-\frac\gamma2v+\Omega_Rw,\qquad
\dot w=-\gamma(w+1)-\Omega_Rv.
\]
稳态有 \(v_{\rm ss}=2\Omega_Rw_{\rm ss}/\gamma\)、
\(w_{\rm ss}=-\gamma^2/(\gamma^2+2\Omega_R^2)\)，所以
\[
p_e^{\rm ss}=\frac{1+w_{\rm ss}}2
=\frac{\Omega_R^2}{\gamma^2+2\Omega_R^2},
\qquad\bar J=\gamma p_e^{\rm ss}.
\]
此模型仅含自发发射通道，未另计探测效率、纯退相干和非共振失谐。

(2) 单原子发出光子后由 \(\sigma_-\rho\sigma_+\) 条件复位为
\(|g\rangle\langle g|\)。于是标准正则光子归一化相关可写成
\[
g^{(2)}(\tau)=\frac{p_e(\tau\mid g)}{p_e^{\rm ss}},\qquad\tau\ge0.
\]
求解上一问的两个Bloch微分方程，初始 \(v(0)=0,w(0)=-1\)。线性化矩阵的两本征值为
\(-3\gamma/4\pm i\mu\)，其中
\(\mu=\sqrt{\Omega_R^2-\gamma^2/16}\)。
故
\[
\boxed{g^{(2)}(\tau)=1-e^{-3\gamma\tau/4}
\left[\cos(\mu\tau)+\frac{3\gamma}{4\mu}\sin(\mu\tau)\right]}.
\]
\(\Omega_R=\gamma/4\) 时用 \(\sin(\mu\tau)/\mu\to\tau\)；
更弱驱动时 \(\mu=i\nu\)，相应变成
\(1-e^{-3\gamma\tau/4}[\cosh(\nu\tau)+3\gamma\sinh(\nu\tau)/(4\nu)]\)，为实数且连续。

(3) 由于探测瞬间激发态占据严格为零，\(g^{(2)}(0)=0\)。
全部相关的指数模均衰减（\(\gamma>0,\Omega_R>0\)），因此
\(\lim_{\tau\to\infty}g^{(2)}(\tau)=1\)。
直接级数展开得到
\[
g^{(2)}(\tau)=\frac{\gamma^2+2\Omega_R^2}{4}\tau^2+O(\tau^3).
\]
短时间二次增长反映原子必须先重新受激再能发第二个光子。
强驱动 \(\Omega_R\gg\gamma\) 时，
\(\mu\simeq\Omega_R\)，\(g^{(2)}\) 可以阻尼振荡并在某些延迟超过1；零延迟严格为0是非经典反聚束，不满足把经典复高斯光场代入后得到的 \(g^{(2)}(0)\ge1\) 结构。

(4) 令 \(a=3\gamma/4\)，则
\(\int_0^\infty e^{-a\tau}\cos\mu\tau\,d\tau=a/(a^2+\mu^2)\)、
\(\int_0^\infty e^{-a\tau}\sin\mu\tau\,d\tau=\mu/(a^2+\mu^2)\)，因此
\[
\int_0^\infty[g^{(2)}(\tau)-1]d\tau
=-\frac{2a}{a^2+\mu^2}
=-\frac{3\gamma}{\gamma^2+2\Omega_R^2}.
\]
代入光子通量 \(\bar J=\gamma\Omega_R^2/(\gamma^2+2\Omega_R^2)\) 可得
\[
\boxed{F=1-\frac{6\gamma^2\Omega_R^2}
{(\gamma^2+2\Omega_R^2)^2}}.
\]
令 \(y=2\Omega_R^2/\gamma^2\)，校正项为 \(3y/(1+y)^2\)，在 \(y=1\) 最大，故
\(F_{\min}=1/4\) 发生于 \(\Omega_R=\gamma/\sqrt2\)。
各衰减指数实部为负，积分严格收敛；亚泊松抑噪来自每次辐射跳跃复位、下一光子之前的重新激发时间，不应与独立Poisson泊松过程的 \(F=1\) 混淆。
